% Systemdokumentation OOP3
% Unterlage für Studenten als Leitfaden für die Erstellung einer SystemDoku
% 17. Oktober 2022
% ---------------------------------------------------------------------------

% --- Templates -------------------------------------------------------------
%\includegraphics{}


% Dokumentklasse
% --------------
\documentclass[12pt,naustrian,a4widepaper]{scrartcl}   
% article style
%   - 11pt Schriftgroesse
%   - new austrian (neue Rechtschreibung)
%   - Papierformat A4
%   - pdf-hyperlinks

% Packages
% --------
%\RequirePackage{hgblistings}
\usepackage[utf8]{inputenc}  % fuer Umlaute, äöü
\usepackage[T1]{fontenc}
\usepackage{a4wide}
\usepackage{times}      % Times Schriften (zusammen mit fontencoding, s.o.)
\usepackage{babel}
\usepackage{graphicx}	  % für das Einbinden von Grafiken
\usepackage{color}      % für färbigen Text
\usepackage{framed}     % für (Text-) Rahmen
\usepackage{listings}   % für den Sourcecode
\usepackage{fancyhdr}   % für Kopf- und Fusszeilen
\usepackage{pdfpages}	% pdf import fuer angabe

\pagestyle{fancy}       % Kopf- / Fusszeile aktivieren
\lstset{
  language={C++},
  basicstyle=\footnotesize\ttfamily,
  keywordstyle=\color{blue},%\bfseries,
  commentstyle=\color{green},
  frame=single,
  xleftmargin=0pt,
  framexleftmargin=0pt,
  framexrightmargin=0pt,
  linewidth=\dimexpr\textwidth+3cm\relax,
  breaklines=true,
  tabsize=3,
  numbers=left, numberstyle=\tiny, stepnumber=1, numbersep=5pt
}

% Seitenspiegel -------------
\typearea{8}	% Festlegung des Seitenspiegels gem. Koma. 4..groß, 9..klein

% Kopfzeile ---------
\lhead{{\footnotesize{Boris Dobrijevic, Maximilian Ebert}}}   %  (links)
\rhead{{\footnotesize{Seite \thepage}}}      %  (rechts)

% Fusszeile ---------
\lfoot{}  % links
\cfoot{}  % mitte 
\rfoot{}  % rechts

% Absatzformatierung ------------------
\setlength{\parindent}{0cm}   % Einrückung der 1. Zeile jedes Absatzes
\setlength{\parskip}{10pt}    % Abstand zwischen den Absätzen


% Package für Hyperlinks (mit pdf-Optionen) -----------------------------------------
\usepackage[
urlcolor=blue,		% blaue weblinks
linkcolor=black,	% interne Links sind schwarz
colorlinks=true,        % links werden eingefürbt
pdfstartview=FitH,      % PDF-Anzeige: Fensterbreite
pdfborder={0 0 0},      % keine Umrandung um links
pdftitle   ={Systemdokumentation},% Referenzen in der Pdf-Datei
pdfauthor  ={Boris Dobrijevic, Maximilian Ebert},
pdfsubject ={Systemdoku},
pdfcreator ={Der Creator},
pdfproducer={Der Producer},
pdfkeywords={Dokumentation, Systemdokumentation}
]{hyperref}

% Beginn des Dokumentes ---------------------
\begin{document}
\selectlanguage{naustrian}   % oder "american" für engl. Texte

% Titelblatt ----------
\title {\vspace{1cm}
       \includegraphics[width=8cm]{./Images/FhOOeLogoOkt2009_HSD_Rot_pastell}\\
       \vspace{2cm}
       {\textbf{Systemdokumentation\\Übung 02}}\\
       \vspace{5mm}
       {\small{Version 1.0}}\\
       \vspace{5mm}
}

\author{\small{Boris Dobrijevic}}
\date  {\small{Hagenberg, \today}}
%\maketitle


\includepdf[pages={1-4}]{uebung_2_RTO.pdf}
\clearpage

% Inhalts-, Tabellen- und Bildverzeichnis (werden generiert) ----------------------------------------------------------
\tableofcontents

\clearpage
\subsection{Schritt 1}
Alle Files wurden wie in der Angabe abverlangt angelegt und die verlangten Funktionen wurden implementiert.

Die Funktion APOS\_TASK\_Create() füllt einfach das spezifizierte TCB Struct mit den übergebenen Werten an.
Mehr zu den TCB Strukturen in Schritt 2.
\\
\\
Die Process-Stacks sind global angelegt und werden mit der Funktion APOS\_Init() initialisiert. Hierbei werden die gesamten Process-Stacks mit 0xFF beschrieben. 

Nach dem Anlegen der TCBs und dem Initialisieren der Stacks kann die Funktion APOS\_Start() ausgeführt werden.
Diese vollführt dabei wie bereits in Übung 1 einen Task-Switch vom Master-Stack auf den Process-Stack, hierbei initial auf TaskA:\\
\begin{lstlisting}[basicstyle=\ttfamily\footnotesize,breaklines=true,frame=single,label=regular_loop]

    __set_PSP((uint32_t)TCB_TaskA.pStack + TCB_TaskA.StackSize);  // PSP zeigt ans Ende des Arrays
    __set_CONTROL(0x02);                    // Bit 1 = 1 ? benutze PSP im Thread-Modus
    __ISB();                                // Pipeline leeren
\end{lstlisting}

Bei den FillTasks war zu beachten, dass die oberen Register nicht einfach direkt beschrieben werden können.
Hierbei werden indirekte Verschiebungen mittels MOVS durchgefürht.\\
\begin{lstlisting}[basicstyle=\ttfamily\footnotesize,breaklines=true,frame=single,label=regular_loop]
void FillTaskC(void)
{
    __asm volatile (
        "movs r0,  #0xC8\n"
        "mov  r8,  r0\n"
        ...
}
\end{lstlisting}


\subsubsection{APOS}
Im APOS-Header-File wurde eine Struktur für TCBs angelegt, welche zum Verwalten der Tasks genutzt werden.
Dabei wurde die Funktionsdeklaration APOS\_TASK\_Create() als Vorlage dafür herbeigezogen.

Die Strukturen der Tasks sind global angelegt und werden bei Bedarf per "extern" Keyword eingebunden.
Somit sind diese immer garantiert angelegt.
\begin{lstlisting}[basicstyle=\ttfamily\footnotesize,breaklines=true,frame=single,label=regular_loop]
typedef struct {
    char const* pTaskName;
    uint32_t Priority;
    void (*pRoutine)(void);
    void * pStack;
    uint32_t StackSize;
    uint32_t TimeSlice;
} APOS_TCB_STRUCT;
\end{lstlisting}





In der Funktion APOS\_Init() 


% #################################################################################################
%      Quellcode
% #################################################################################################
\clearpage
\section{Quellcode}
% UnitTest
\subsection{Beispiel 1}
% \subsubsection{UnitTest - main.cpp}
% \lstinputlisting{./../UnitTest/main.cpp}
% 
% \subsubsection{UnitTest - StringOps.h}
% \lstinputlisting{./../UnitTest/StringOps.h}
% 
% \subsubsection{UnitTest - StringOps.cpp}
% \lstinputlisting{./../UnitTest/StringOps.cpp}
% 
% \subsubsection{UnitTest - StringOpsTest.h}
% \lstinputlisting{./../UnitTest/StringOpsTest.h}
% 
% \subsubsection{UnitTest - StringOpsTest.cpp}
% \lstinputlisting{./../UnitTest/StringOpsTest.cpp}
%       
% 
% \subsection{Beispiel 2}
% \subsubsection{Laufzeit - main.cpp}
% \lstinputlisting{./../Laufzeit/main.cpp}
% 
% \subsubsection{Laufzeit - Laufzeit.h}
% \lstinputlisting{./../Laufzeit/Laufzeit.h}
% 
% \subsubsection{Laufzeit - Laufzeit.cpp}
% \lstinputlisting{./../Laufzeit/Laufzeit.cpp}
% 
% 
% \subsection{Beispiel 3}
% \subsubsection{ListVsVector - main.cpp}
% \lstinputlisting{./../ListVsVector/main.cpp}
% 
% \subsubsection{ListVsVector - ListVsVector.h}
% \lstinputlisting{./../ListVsVector/ListVsVector.h}
% 
% \subsubsection{ListVsVector - ListVsVector.h}
% \lstinputlisting{./../ListVsVector/ListVsVector.cpp}

\end{document}
